\section{Conclusion}
\label{sec:conclusion}

We introduced \method, a cheap, controllable counterfactual testbed for evaluating vision-language models. By pairing each synthetic scene with a minimally intervened version and measuring whether models change their answers \emph{if and only if} the ground truth changes, \method provides a fundamentally different evaluation axis from standard accuracy metrics. Our Counterfactual Consistency Score (\ccs), decomposed into sensitivity and specificity, reveals systematic failures that accuracy alone misses: a 19.5-point \ccs gap across six VLMs, with two distinct failure profiles---Claude~3.5 Sonnet suffers from low specificity (spurious changes) while Llama~3.2 11B suffers from low sensitivity (sticky answers). Notably, the open-weight Qwen2.5-VL~72B rivals proprietary GPT-4o, demonstrating that scale and training recipe, not proprietary status, determine counterfactual consistency.

\paragraph{Limitations.}
\method uses simple 2D synthetic scenes that lack the visual complexity of real-world images. Performance on \method may not directly predict performance on natural images. The benchmark currently evaluates only binary (yes/no) and numerical (counting) questions; extending to free-form answers would require more sophisticated evaluation. While we evaluate six models spanning proprietary and open-weight families, future work should cover an even broader range of architectures and scales.

\paragraph{Future Work.}
Several extensions are natural: (i) adding difficulty levels by increasing scene complexity (more objects, smaller sizes, confusable colors), (ii) extending to 3D rendered scenes for more realistic visual reasoning, (iii) testing multi-step counterfactual chains (``if A were removed, and then B were moved, would C still be visible?''), and (iv) using \method-style evaluation as a training signal for improving counterfactual consistency. We also plan to extend \method to multilingual settings and to test whether fine-tuning on counterfactual pairs improves consistency on natural images.

\paragraph{Broader Impact.}
\method contributes to the growing toolkit for evaluating VLM reliability. As these models are deployed in safety-critical applications (medical imaging, autonomous driving, accessibility), understanding the conditions under which their outputs are consistent is essential. We release all code, data, and evaluation scripts to facilitate reproducibility and extension.
